\documentclass{article}

\usepackage{makecell}

\begin{document}

\section*{Ch. 14 HW \#2 Pt. 2}

\subsection*{Instructions} This is not the entirety of Ch. 14 HW \#2. This is just a portion of the assignment that we started in class.

\subsection*{Quick Quiz}

\begin{enumerate}

\item A firm is producing 1,000 units at a total cost of \$5,000. When it increases production to 1,001 units, its total cost rises to \$5,008. For this firm,

	\begin{enumerate}
	
	\item marginal cost is \$5, and average variable cost is \$8.
	
	\item marginal cost is \$8, and average variable cost is \$5. 
	
	\item marginal cost is \$5, and average total cost is \$8.
	
	\item marginal cost is \$8, and average total cost is \$5.
	
	\end{enumerate}
	
\item A firm is producing 20 units with an average total cost of \$25 and a marginal cost of \$15. If it increases production to 21 units, which of the following must occur?

	\begin{enumerate}
	
	\item Marginal cost will decrease.
	
	\item Marginal cost will increase.
	
	\item Average total cost will decrease.
	
	\item Average total cost will increase.
	
	\end{enumerate}
	
\item The government imposes a \$1,000 per year license fee on all pizza restaurants. As a result, which cost curves shift?

	\begin{enumerate}
	
	\item average total cost and marginal cost
	
	\item average total cost and average fixed cost
	
	\item average variable cost and marginal cost
	
	\item average variable cost and average fixed cost
	
	\end{enumerate}

\end{enumerate}

\subsection*{Problems and Applications}

\begin{enumerate}

\item This chapter discusses many types of costs: opportunity cost, total cost, fixed cost, variable cost, average total cost, and marginal cost. Fill in the type of cost that best completes each sentence:

	\begin{enumerate}
	
	\item What you give up in taking some action is called the \underline{\phantom{xxxxxxxxxx}}.
	
	\item \underline{\phantom{xxxxxxxxxx}} is falling when marginal cost is below it and rising when marginal cost is above it.
	
	\item A cost that does not depend on the quantity produced is a(n) \underline{\phantom{xxxxxxxxxx}}.
	
	\item In the ice-cream industry in the short run, \underline{\phantom{xxxxxxxxxx}} includes the cost of cream and sugar but not the cost of the factory.
	
	\item Profits equal total revenue minus \underline{\phantom{xxxxxxxxxx}}.
	
	\item The cost of producing an extra unit of output is the \underline{\phantom{xxxxxxxxxx}}.
	
	\end{enumerate}
	
\item Nimbus, Inc. makes brooms and sells them door-to-door. Here is the relationship between the number of workers and Nimbus's output during a given day:

	\begin{center}
	\begin{tabular}{| c | c | c | c | c | c |}
	\hline
	Workers & Output & \makecell{Marginal\\Product} & \makecell{Total\\Cost} & \makecell{Average\\Total\\Cost} & \makecell{Marginal\\Cost} \\
	\hline
	0 & 0 & & & & \\
	\hline
	1 & 20 & & & &\\
	\hline
	2 & 50 & & & &\\
	\hline
	3 & 90 & & & &\\
	\hline
	4 & 120 & & & &\\
	\hline
	5 & 140 & & & &\\
	\hline
	6 & 150 & & & &\\
	\hline
	7 & 155 & & & & \\
	\hline
	\end{tabular}
	\end{center}
	
	\begin{enumerate}
	
	\item Fill in the column of marginal products. What pattern do you see? How might you explain it?
	
	\item A worker costs \$100 a day, and the firm has fixed costs of \$200. Use this information to fill in the column for total cost.
	
	\item Fill in the column for average total cost. What pattern do you see?
	
	\item Now fill in the column for marginal cost. What pattern do you see?
	
	\item Compare the column for marginal product with the column for marginal cost. Explain the relationship.
	
	\item Compare the column for average total cost with the column for marginal cost. Explain the relationship.
	
	\end{enumerate}
	
\item You are the chief financial officer for a firm that sells gaming consoles. Your firm has the following average-total-cost schedule:

	\begin{center}
	\begin{tabular}{ | c | c |}
	\hline
	Quantity & Average Total Cost\\
	\hline
	600 consoles & \$300 \\
	\hline
	601 & 301\\
	\hline
	\end{tabular}
	\end{center}
	
	Your current level of production is 600 consoles, all of which have been sold. Someone calls, desperate to buy one of your consoles. The caller offers you \$550 for it. Should you accept the offer? Why or why not?

\end{enumerate}

\end{document}